% Part: proof-theory
% Chapter: natural-deduction
% Section: rules-N1

\documentclass[../../../include/open-logic-section]{subfiles}

\begin{document}

\begin{table}
\begin{center}
\begin{tabular}[t]{@{}cc@{}}
\bottomAlignProof
\AxiomC{$\Discharge{!A}{\OLId{latin-ordinary}{x}}$}
\shortDeduce
\DeduceC{$!B$}
\DischargeRule{\Intro{\lif}}{x}
\UnaryInfC{$!A \lif !B$}
\DisplayProof
&
\bottomAlignProof
\AxiomC{$!A \lif !B$}
\AxiomC{$!A$}
\RightLabel{\Elim{\lif}}
\BinaryInfC{$!B$}
\DisplayProof
\\[4ex]
\bottomAlignProof
\AxiomC{$!A$}
\AxiomC{$!B$}
\RightLabel{\Intro{\land}}
\BinaryInfC{$!A \land !B$}
\DisplayProof
&
\begin{tabular}[b]{@{}l@{}}
\AxiomC{$!A \land !B$}
\RightLabel{\Elim{\land}[1]}
\UnaryInfC{$!A$}
\DisplayProof
\\[3ex]
\bottomAlignProof
\AxiomC{$!A \land !B$}
\RightLabel{\Elim{\land}[2]}
\UnaryInfC{$!B$}
\DisplayProof
\end{tabular}
\\[4ex]
\begin{tabular}{@{}l@{}}
\AxiomC{$!A$}
\RightLabel{\Intro{\lor}[1]}
\UnaryInfC{$!A \lor !B$}
\DisplayProof
\\[3ex]
\AxiomC{$!B$}
\RightLabel{\Intro{\lor}[2]}
\UnaryInfC{$!A \lor !B$}
\DisplayProof
\end{tabular}
&
\AxiomC{$!A \lor !B$}
\AxiomC{$\Discharge{!A}{\OLId{latin-ordinary}{x}}$}
\shortDeduce
\DeduceC{$!C$}
\AxiomC{$\Discharge{!B}{\OLId{latin-ordinary}{y}}$}
\shortDeduce
\DeduceC{$!C$}
\DischargeRule{\Elim{\lor}}{x\,y}
\TrinaryInfC{$!C$}
\DisplayProof
\bottomAlignProof
\\[6ex]
\bottomAlignProof
\AxiomC{$\Discharge{!A}{\OLId{latin-ordinary}{x}}$}
\shortDeduce
\DeduceC{$\lfalse$}
\DischargeRule{\Intro{\lnot}}{x}
\UnaryInfC{$\lnot !A$}
\DisplayProof
&
\bottomAlignProof
\AxiomC{$\lnot !A$}
\AxiomC{$!A$}
\RightLabel{\Elim{\lnot}}
\BinaryInfC{$\lfalse$}
\DisplayProof
\\[4ex]
\bottomAlignProof
\AxiomC{$!A(\OLId{latin-ordinary}{a})$}
\RightLabel{\Intro{\lforall}}
\UnaryInfC{$\lforall[\OLId{latin-ordinary}{x}][!A(\OLId{latin-ordinary}{x})]$}
\DisplayProof
&
\bottomAlignProof
\AxiomC{$\lforall[\OLId{latin-ordinary}{x}][!A(\OLId{latin-ordinary}{x})]$}
\RightLabel{\Elim{\lforall}}
\UnaryInfC{$!A(\OLId{latin-ordinary}{t})$}
\DisplayProof
\\[4ex]
\bottomAlignProof
\AxiomC{$!A(\OLId{latin-ordinary}{t})$}
\RightLabel{\Intro{\lexists}}
\UnaryInfC{$\lexists[\OLId{latin-ordinary}{x}][!A(\OLId{latin-ordinary}{x})]$}
\DisplayProof
&
\bottomAlignProof
\AxiomC{$\lexists[\OLId{latin-ordinary}{x}][!A(\OLId{latin-ordinary}{x})]$}
\AxiomC{$\Discharge{!A(\OLId{latin-ordinary}{a})}{\OLId{latin-ordinary}{x}}$}
\shortDeduce
\DeduceC{$!C$}
\DischargeRule{\Elim{\lexists}}{x}
\BinaryInfC{$!C$}
\DisplayProof
\\[4ex]
\bottomAlignProof
\AxiomC{$\lfalse$}
\RightLabel{\FalseInt}
\UnaryInfC{$!A$}
\DisplayProof
&
\bottomAlignProof
\AxiomC{$\Discharge{\lnot!A}{\OLId{latin-ordinary}{x}}$}
\shortDeduce
\DeduceC{$\lfalse$}
\DischargeRule{\FalseCl}{x}
\UnaryInfC{$!A$}
\DisplayProof
\end{tabular}
\end{center}
\caption{هذه قواعد~\Log{N1i} و\Log{N1c}. ويشترط في~$\Intro{\lforall}$
ألا يظهر~!!{constant}~$\OLId{latin-ordinary}{a}$ في النتيجة ولا في فرض مفتوح. وأما
$\Elim{\lexists}$ فيشترط فيه ألا يظهر~$\OLId{latin-ordinary}{a}$ في المقدمة الكبرى ولا في
النتيجة، ولا في فرض مفتوح سوى الفرض~$!A(\OLId{latin-ordinary}{a})^\OLId{latin-ordinary}{x}$ الذي يبرئه الاستدلال.
وليس في~\Log{N1i} القاعدة~\FalseCl.}
\ollabel{tab:N1}
\end{table}

\end{document}
